\section{Conditional moments and the one-sided comparison}
\label{app:conditional}
This appendix proves the estimates used in
Section~\ref{sec:local}. All conditioning is on the actual Bernoulli
count. Dependence between a label and its residual is preserved.

\subsection{Exponential tilting and projection}
Retain $p_i\in[1/4,3/4]$, $v_i=p_i(1-p_i)$,
$\sum_i v_i a_i=0$, and $\Lambda_A=\sum_i v_i a_i^2$.
Suppose $m\ge10^{12}$ and $|k-\mu|\le6\sqrt m$.
Tilt by $e^{\theta K_0}$ so that the count mean is $k$, putting
\[
 r_i=\frac{p_i e^\theta}{1-p_i+p_i e^\theta},\quad
 w_i=r_i(1-r_i),\quad V_\theta=\sum_i w_i.
\]
For $|\theta|\le\log2$, $r_i\in[1/8,7/8]$ and $w_i\ge7/64>1/10$.
The count mean changes by more than $6\sqrt m$ in both directions
before that boundary, so the tilt exists and
\begin{equation}\label{eq:tilt-quantities}
 |\theta|\le60/\sqrt m,\quad |w_i/v_i-1|\le2|\theta|,
       \quad |r_i-p_i-\theta v_i|\le\theta^2v_i.
\end{equation}
The conditional law given $K_0=k$ is unchanged. Define
\[
 c_\theta=\frac{\sum_iw_i a_i}{V_\theta},\quad b_i=a_i-c_\theta,
 \quad L=\sum_i b_i(B_i-r_i),\quad Q=\sum_i b_i^2,
 \quad Q_\theta=\sum_iw_i b_i^2.
\]
On the count event, $A=m_k^0+L$, where
\begin{equation}\label{eq:tilted-gap-data}
 |m_k^0|=\left|\sum_i a_i(r_i-p_i)\right|
      \le1800\sqrt{\Lambda_A}/\sqrt m,\quad
 |Q_\theta-\Lambda_A|\le180\Lambda_A/\sqrt m,
 \quad Q\le16\Lambda_A.
\end{equation}
For the second bound expand
$Q_\theta=\sum_iw_i a_i^2-(\sum_iw_i a_i)^2/V_\theta$.
The errors are at most $2|\theta|\Lambda_A$ and
$8\theta^2\Lambda_A$, together at most $3|\theta|\Lambda_A$.
Also $\Lambda_A/2\le Q_\theta\le3\Lambda_A/2$; using
$w_i\ge7/64$ proves the last bound. If $\Lambda_A=0$, all the
gap estimates are immediate, so suppose it is positive.

\subsection{Fourier moments with their residual factors}
For $|t|\le1/4$, put
$z_i(t)=r_ie^{it}/(1-r_i+r_ie^{it})$.
The first centered complex cumulant is $c_{i,1}=z_i-r_i$; the next
three are
\[
 c_{i,2}=z_i(1-z_i),\quad
 c_{i,3}=z_i(1-z_i)(1-2z_i),\quad
 c_{i,4}=z_i(1-z_i)(1-6z_i(1-z_i)).
\]
Both $z_i$ and $1-z_i$ have modulus at most two. Differentiation gives
\[
 |c_{i,1}-iw_it|\le10t^2,\quad |c_{i,2}-w_i|\le20|t|,
 \quad |c_{i,3}|\le20,\quad |c_{i,4}|\le100.
\]
For $M_j=\sum_i b_i^jc_{i,j}$, the exact cancellation
$\sum_iw_i b_i=0$ yields
\begin{gather}
 |M_1|\le10\sqrt{mQ}t^2,\quad
 |M_2-Q_\theta|\le20Q|t|,\quad |M_2|\le6Q,\notag\\
 |M_3|\le20Q^{3/2},\qquad |M_4|\le100Q^2.              \label{eq:complex-moments}
\end{gather}
The tilted count transform $F_\theta$ satisfies
$|F_\theta(t)|\le e^{-mt^2/50}$ on $[-\pi,\pi]$.
Lemma~\ref{lem:count-estimate} gives
$\Pp_\theta(K_0=k)\ge1/(3\sqrt m)$.
After division by this probability, upper bounds for the Gaussian
integrals of powers $0,1,2,4,8$ are respectively
\begin{equation}\label{eq:gaussian-moment-table}
             7,\quad25m^{-1/2},\quad175m^{-1},\quad
                          13125m^{-2},\quad3\cdot10^8m^{-4}.
\end{equation}
These follow by scaling $t\sqrt m$ in
$(2\pi)^{-1}\int e^{-mt^2/50}|t|^jdt$.
On $1/4\le|t|\le\pi$, expand the moment before factoring the
product. A term of degree $j\le4$ removes at most four Bernoulli
factors; the remaining product is at most $e^{-m/810}$.
The conditional tail contribution is therefore at most
$3m^{(j+1)/2}Q^{j/2}e^{-m/810}$.
For $j=1,2,4$ and $m\ge10^8$, these are at most
$\sqrt Q/\sqrt m$, $Q/\sqrt m$, and $Q^2$.
For instance $e^{-m/810}\le10^{18}/m^5$ reduces the verification
to decreasing powers of $m$.

The first, second, and fourth moment numerators are exactly
\[
 F_\theta M_1,\quad F_\theta(M_1^2+M_2),\quad
 F_\theta(M_1^4+6M_1^2M_2+3M_2^2+4M_1M_3+M_4).
\]
Equations \eqref{eq:complex-moments}--\eqref{eq:gaussian-moment-table},
with the tails just bounded, imply
\begin{equation}\label{eq:projected-moments}
 \begin{gathered}
 |\E(L\mid K_0=k)|\le2000\sqrt Q/\sqrt m,\\
 |\Var(L\mid K_0=k)-Q_\theta|\le2000Q/\sqrt m,
 \qquad \E(L^4\mid K_0=k)\le2000Q^2.
 \end{gathered}
\end{equation}
Here is numerical detail for these enlargements.
The mean's central coefficient is $1750$.
The second moment errors are at most
$500Q/\sqrt m+1312500Q/m+2Q/\sqrt m$,
and the squared mean costs at most $4\cdot10^6Q/m$.
The five fourth-moment central coefficients, in the displayed order,
are at most
\[
 3\cdot10^{12}/m^2,\quad47250000/m,\quad756,\quad
                         140000/\sqrt m,\quad700,
\]
all multiplying $Q^2$. Their sum and the tail are below $2000Q^2$.
Centering raises the fourth-moment bound to $20000Q^2$.
Combining \eqref{eq:tilted-gap-data} and
\eqref{eq:projected-moments} proves
\eqref{eq:conditional-mean}--\eqref{eq:conditional-gap}.

For completeness, the relative label estimate
\eqref{eq:conditional-label} also follows numerically.
Under the tilt, deleting a Bernoulli changes count variance by at most
$1/4$ and displaces the new threshold from its mean by at most one.
The Gaussian mass ratio has logarithm bounded by $1/V_\theta$.
The $6/V$ errors and a lower bound by half the centered Gaussian mass
then give a deleted-count ratio error at most
\[
 4/V_\theta+28/(\phz\sqrt{V_\theta})\le300/\sqrt m.
\]
Each tilted label probability differs relatively from its original
probability by at most $2|\theta|\le120/\sqrt m$.
Multiplication gives $500/\sqrt m$ for $m\ge10^{12}$.
This is a relative bound, so arbitrarily small within-cluster noise
variances remain controlled after conditioning.

\subsection{Proof of the one-sided loss}
We prove \eqref{eq:one-sided}, following the clipping method of
\citet[Lemma~3.2 and equation~(83)]{HeCheng2026}.
First suppose $V_0$ is centered and $|V_0|\le a=1/10$.
For $u\ge0$, $\E(V_0)_+\le u+a\Pp(V_0>u)$ and
$f(u)\ge3u/4$ give
$\Pp(V_0>u)+f(u)\ge3\E|V_0|/8$.
If $u=-r$, $0<r\le a$, put $t=\Pp(V_0>-r)$.
Centering gives $0\le at-r(1-t)$, hence $r\le2at$.
Also $\E|V_0|=2\E(V_0)_+\le2at$, so
$t+f(-r)\ge t-5r/4\ge3t/4\ge3\E|V_0|/8$.
For $a<r\le13/20$, the probability is one and
$1+f(-r)\ge3/16\ge3\E|V_0|/8$.

For general centered $Y$, clip at $1/20$, obtaining $T$, and put
$b=\E T$, $V_0=T-b$, $M_4=\E Y^4$.
Then
\[
 |b|\le8000M_4,\quad\Pp(T\ne Y)\le160000M_4,\quad
                  \E|V_0|\ge\E|Y|-16000M_4.
\]
If $M_4\le1/160000$, then $|b|\le1/20$,
$|V_0|\le1/10$, and $|u-b|\le13/20$.
Apply the bounded case at $u-b$, and compare $f(u-b)$ with $f(u)$.
The total error coefficient is
$(3/8)16000+160000+(5/4)8000=176000$.
This proves the first inequality in \eqref{eq:one-sided}.
Apply it to $-Y$ and $-f(-\,\cdot\,)$ for the second.
All events in this proof are strict; atoms at the threshold are included
with the conventions needed in \eqref{eq:layers}.
